% ========================================================
% BÀI TẬP TOÁN 5 - KẾT NỐI TRI THỨC VỚI CUỘC SỐNG
% Bài 63: Thu thập, phân loại, sắp xếp các số liệu
% Định dạng: BTTeX Part III - Trắc nghiệm điền số ngắn
% ========================================================

% Câu 1
\begin{ex}
  Thống kê điểm kiểm tra Toán của 10 học sinh: 8, 9, 7, 10, 8, 9, 6, 8, 9, 10. Có bao nhiêu học sinh đạt điểm 10?
  \shortans{2}
  \loigiai{
    Có hai điểm 10.
  }
\end{ex}

% Câu 2
\begin{ex}
  Ở bảng thống kê trên, có bao nhiêu học sinh đạt điểm 8?
  \shortans{3}
  \loigiai{
    Có ba bạn điểm 8.
  }
\end{ex}

% Câu 3
\begin{ex}
  Ở bảng thống kê trên, điểm số thấp nhất là bao nhiêu điểm?
  \shortans{6}
  \loigiai{
    Điểm thấp nhất là 6.
  }
\end{ex}

% Câu 4
\begin{ex}
  Ở bảng thống kê trên, hiệu giữa điểm cao nhất và thấp nhất là bao nhiêu?
  \shortans{4}
  \loigiai{
    $10 - 6 = 4$.
  }
\end{ex}

% Câu 5
\begin{ex}
  Ở bảng trên, có bao nhiêu bạn đạt điểm 9?
  \shortans{3}
  \loigiai{
    Có 3 bạn đạt điểm 9.
  }
\end{ex}

% Câu 6
\begin{ex}
  Số cây trồng của 4 tổ lần lượt là: Tổ 1: 35 cây, Tổ 2: 40 cây, Tổ 3: 45 cây, Tổ 4: 40 cây. Cả bốn tổ trồng được bao nhiêu cây?
  \shortans{160}
  \loigiai{
    $35 + 40 + 45 + 40 = 160$ cây.
  }
\end{ex}

% Câu 7
\begin{ex}
  Ở bài trên, trung bình mỗi tổ trồng được bao nhiêu cây?
  \shortans{40}
  \loigiai{
    $160 : 4 = 40$ cây.
  }
\end{ex}

% Câu 8
\begin{ex}
  Ở bài trên, tổ trồng được nhiều cây nhất trồng hơn tổ trồng ít nhất bao nhiêu cây?
  \shortans{10}
  \loigiai{
    $45 - 35 = 10$ cây.
  }
\end{ex}

% Câu 9
\begin{ex}
  Thống kê số học sinh yêu thích các môn học của lớp 5A: Toán: 15 bạn, Tiếng Việt: 10 bạn, Tiếng Anh: 12 bạn, Tin học: 8 bạn. Môn học được yêu thích nhất có bao nhiêu bạn chọn?
  \shortans{15}
  \loigiai{
    Toán có 15 bạn chọn (nhiều nhất).
  }
\end{ex}

% Câu 10
\begin{ex}
  Ở thống kê trên, tổng số học sinh thích môn Toán và Tin học là bao nhiêu bạn?
  \shortans{23}
  \loigiai{
    $15 + 8 = 23$ bạn.
  }
\end{ex}

% Câu 11
\begin{ex}
  Khối lớp 5 có 4 lớp quyên góp số sách: 5A: 120 quyển, 5B: 110 quyển, 5C: 130 quyển, 5D: 140 quyển. Tổng số sách quyên góp được là bao nhiêu quyển?
  \shortans{500}
  \loigiai{
    $120 + 110 + 130 + 140 = 500$ quyển.
  }
\end{ex}

% Câu 12
\begin{ex}
  Ở bài trên, trung bình mỗi lớp quyên góp được bao nhiêu quyển sách?
  \shortans{125}
  \loigiai{
    $500 : 4 = 125$ quyển.
  }
\end{ex}

% Câu 13
\begin{ex}
  Chiều cao của 5 bạn học sinh (đơn vị cm) là: 138, 142, 140, 145, 135. Chiều cao của bạn cao nhất là bao nhiêu cm?
  \shortans{145}
  \loigiai{
    Bạn cao nhất là 145 cm.
  }
\end{ex}

% Câu 14
\begin{ex}
  Ở bài trên, chiều cao trung bình của 5 bạn là bao nhiêu cm?
  \shortans{140}
  \loigiai{
    $(138 + 142 + 140 + 145 + 135) : 5 = 700 : 5 = 140\text{ cm}$.
  }
\end{ex}

% Câu 15
\begin{ex}
  Số bạn nam của 3 tổ lần lượt là 6, 7, 8 bạn. Số bạn nữ của 3 tổ lần lượt là 7, 6, 6 bạn. Tổng số học sinh của cả 3 tổ là bao nhiêu bạn?
  \shortans{40}
  \loigiai{
    $(6+7+8) + (7+6+6) = 21 + 19 = 40$ bạn.
  }
\end{ex}

% Câu 16
\begin{ex}
  Thống kê số huy chương của một đoàn thể thao: Vàng: 15, Bạc: 12, Đồng: 18. Tổng số huy chương đoàn giành được là bao nhiêu?
  \shortans{45}
  \loigiai{
    $15 + 12 + 18 = 45$.
  }
\end{ex}

% Câu 17
\begin{ex}
  Ở bảng trên, số huy chương đồng nhiều hơn huy chương bạc là bao nhiêu chiếc?
  \shortans{6}
  \loigiai{
    $18 - 12 = 6$ chiếc.
  }
\end{ex}

% Câu 18
\begin{ex}
  Một tổ công nhân có mức sản xuất trong 4 tuần (sản phẩm): 240, 260, 250, 250. Tổng số sản phẩm làm được là bao nhiêu?
  \shortans{1000}
  \loigiai{
    $240 + 260 + 250 + 250 = 1\,000$.
  }
\end{ex}

% Câu 19
\begin{ex}
  Ở bài trên, trung bình mỗi tuần tổ đó sản xuất được bao nhiêu sản phẩm?
  \shortans{250}
  \loigiai{
    $1\,000 : 4 = 250$ sản phẩm.
  }
\end{ex}

% Câu 20
\begin{ex}
  Điểm kiểm tra của An qua 5 bài Toán là: 8, 9, 8, 10, 10. Điểm trung bình môn Toán của An là bao nhiêu điểm (làm tròn số nguyên)?
  \shortans{9}
  \loigiai{
    $(8 + 9 + 8 + 10 + 10) : 5 = 45 : 5 = 9$ điểm.
  }
\end{ex}

% Câu 21
\begin{ex}
  Trong đợt tiêm chủng có 3 nhóm: Nhóm A có 25 người, nhóm B có 35 người, nhóm C có 40 người. Số người nhóm C chiếm bao nhiêu phần trăm tổng số người?
  \shortans{40}
  \loigiai{
    Tổng: 100 người. Nhóm C chiếm $40\%$.
  }
\end{ex}

% Câu 22
\begin{ex}
  Ở bảng trên, số người nhóm A chiếm bao nhiêu phần trăm?
  \shortans{25}
  \loigiai{
    $25\%$.
  }
\end{ex}

% Câu 23
\begin{ex}
  Số lượt mượn sách thư viện trong tuần: Thứ Hai: 45 lượt, Thứ Ba: 50 lượt, Thứ Tư: 55 lượt, Thứ Năm: 60 lượt, Thứ Sáu: 40 lượt. Cả tuần có bao nhiêu lượt mượn sách?
  \shortans{250}
  \loigiai{
    $45 + 50 + 55 + 60 + 40 = 250$ lượt.
  }
\end{ex}

% Câu 24
\begin{ex}
  Ở bảng trên, ngày có lượt mượn nhiều nhất hơn ngày ít nhất bao nhiêu lượt?
  \shortans{20}
  \loigiai{
    $60 - 40 = 20$ lượt.
  }
\end{ex}

% Câu 25
\begin{ex}
  Sắp xếp 4 số liệu sau theo thứ tự tăng dần: $15, 8, 23, 12$. Giá trị của số đứng ở vị trí thứ ba là bao nhiêu?
  \shortans{15}
  \loigiai{
    Dãy tăng dần: 8, 12, 15, 23. Số thứ ba là 15.
  }
\end{ex}
